% ---
% 4 Software
% Este arquivo apresenta o resultado prático do projeto: o software implementado.
% Conforme as Diretrizes do PPC, deve incluir: implantação e testes.
% ---
\chapter{SOFTWARE}
\label{cap:software}

Neste capítulo, o autor deve apresentar o resultado prático do projeto: o software implementado. As Diretrizes do PPC estabelecem que esta seção deve demonstrar o funcionamento do sistema, destacando suas principais funcionalidades através de capturas de tela (\textit{screenshots}) e trechos de código relevantes.

\section{Implantação}

Descreva o processo de implantação (\textit{deploy}) do software. Indique o ambiente de hospedagem utilizado (servidor local, nuvem, etc.), as configurações necessárias e os passos para colocar o sistema em produção.

Apresente as telas do sistema em funcionamento. Diferente da prototipagem (Capítulo \ref{cap:modelagem}), aqui devem ser mostradas as interfaces reais do software finalizado.

\begin{figure}[htb]
\centering
\caption{Tela Inicial do Sistema}
\label{fig:tela_inicial}
% \includegraphics[width=0.9\textwidth]{figuras/tela_inicial.png}
\fbox{\parbox{0.9\textwidth}{\centering \vspace{3cm} [Inserir Print da Tela do Sistema aqui] \\ Salve a imagem em: \texttt{figuras/tela\_inicial.png} \vspace{3cm}}}
\source{Elaborada pelo autor (\imprimirdata).}
\end{figure}

% --- EXEMPLO DE CÓDIGO-FONTE ---
% Utilize o ambiente lstlisting para inserir trechos de código.
% O pacote listings já está configurado em configuracoes.tex.
O Código \ref{lst:codigo_auth} apresenta a função responsável pela autenticação de usuários no sistema.

\begin{lstlisting}[language=Python, caption={Função de autenticação de usuários}, label={lst:codigo_auth}]
from django.contrib.auth.hashers import check_password
from .models import Usuario

def autenticar_usuario(username, password):
    """
    Verifica as credenciais do usuario no banco de dados.
    Retorna o objeto Usuario se autenticado, None caso contrario.
    """
    try:
        usuario = Usuario.objects.get(username=username)
        if check_password(password, usuario.password):
            return usuario
    except Usuario.DoesNotExist:
        pass
    return None
\end{lstlisting}

Explique o funcionamento do código apresentado e como ele se integra ao restante da aplicação. Evite colar trechos excessivamente longos ou gerados automaticamente (como \textit{migrations} ou configurações padrão de \textit{frameworks}); foque na lógica de negócios implementada pelo autor.

% --- EXEMPLO DE ALGORITMO (pseudocódigo) ---
% Para algoritmos em pseudocódigo, utilize o ambiente lstlisting
% com language=none ou crie um estilo customizado.
O Algoritmo \ref{lst:algoritmo_busca} apresenta o pseudocódigo do módulo de busca implementado.

\begin{lstlisting}[caption={Algoritmo de busca por palavra-chave}, label={lst:algoritmo_busca}, language={}]
ALGORITMO BuscaPorPalavraChave
ENTRADA: termo (string), registros (lista)
SAIDA: resultados (lista)

INICIO
    resultados <- lista vazia
    PARA CADA registro EM registros FACA
        SE termo ESTA_CONTIDO_EM registro.titulo OU
           termo ESTA_CONTIDO_EM registro.descricao ENTAO
            ADICIONAR registro EM resultados
        FIM_SE
    FIM_PARA
    ORDENAR resultados POR relevancia DECRESCENTE
    RETORNAR resultados
FIM
\end{lstlisting}

\section{Testes}

Descreva os testes realizados para validar o software. Indique o tipo de teste realizado (unitário, integração, sistema, aceitação, usabilidade) e apresente os resultados obtidos.

\begin{table}[htb]
\centering
\caption{Resultados dos testes realizados}
\label{tab:testes}
\begin{tabular}{@{}llll@{}}
\toprule
\textbf{Código} & \textbf{Tipo} & \textbf{Descrição} & \textbf{Resultado} \\ \midrule
T01 & Unitário & Validação do cadastro de usuário & Aprovado \\
T02 & Unitário & Autenticação com credenciais válidas & Aprovado \\
T03 & Integração & Fluxo completo de login e acesso & Aprovado \\
T04 & Usabilidade & Navegação por usuários leigos & Aprovado \\ \bottomrule
\end{tabular}
\source{Elaborada pelo autor (\imprimirdata).}
\end{table}
